\subsection{关于一个元素的局部化}

设 A 是一个环，M 是一个 A-模。对每个元素 \(f \in \mathbf{A}\)，我们将记 \(\mathbf{A}_f = \mathbf{A}[f^{-1}]\)、\(\mathbf{M}_f = \mathbf{M}[f^{-1}] = \mathbf{M} \otimes_{\mathbf{A}} \mathbf{A}[f^{-1}]\)（§2 第 1 与第 2 节）；若 \(\mathbf{S}_f\) 是诸 \(f^n\)（对 \(n \geqslant 0\)）组成之集，则 \(\mathbf{A}_f = \mathbf{S}_f^{-1}\mathbf{A}\)、\(\mathbf{M}_f = \mathbf{S}_f^{-1}\mathbf{M}\)。若 \(f\) 在 \(\mathbf{A}\) 中可逆，则 \(\mathbf{A}_f\)（相应地，\(\mathbf{M}_f\)）典范地等同于 \(\mathbf{A}\)（相应地，\(\mathbf{M}\)）；若 \(f\) 是幂零的，则 \(\mathbf{A}_f = 0\) 且 \(\mathbf{M}_f = 0\)。对每个 A-模同态 \(u: \mathbf{M} \to \mathbf{N}\)，我们记 \(u_f = u \otimes 1: \mathbf{M}_f \to \mathbf{N}_f\)。

设 \(g\) 是 \(\mathbf{A}\) 的另一个元素；\(\mathbf{A}_{fg}\)（相应地，\(\mathbf{M}_{fg}\)）典范地等同于 \((\mathbf{A}_f)_{g/1}\)（相应地，\((\mathbf{M}_f)_{g/1}\)），其中 \(g/1\) 是 \(g\) 在 \(\mathbf{A}_f\) 中的像，且 \(u_{fg}\) 等同于 \((u_f)_{g/1}\)（§2 第 3 节命题 7）。

命题 1. 设 \(f\) 是环 \(\mathbf{A}\) 的一个元素，\(\phi: \mathbf{A} \to \mathbf{A}_f\) 是典范映射。映射 \(^a\phi: \operatorname{Spec}(\mathbf{A}_f) \to \operatorname{Spec}(\mathbf{A})\) 是从 \(\operatorname{Spec}(\mathbf{A}_f)\) 到开子空间 \(\mathbf{X}_f\)（\(\mathbf{X} = \operatorname{Spec}(\mathbf{A})\) 的）上的一个同胚（§4 第 3 节）。

这是 §4 第 3 节命题 13 的推论的一个特例。

命题 2. 设 A 是一个环，\(u: \mathbf{M} \to \mathbf{N}\) 是一个 A-模同态，\(\mathfrak{p}\) 是 A 的一个素理想。

\begin{enumerate}
\def\labelenumi{(\roman{enumi})}
\item
  设 \(u_{\mathfrak{p}} \colon \mathbf{M}_{\mathfrak{p}} \to \mathbf{N}_{\mathfrak{p}}\) 是满射且 \(\mathbf{N}\) 是有限生成的。则存在 \(f \in \mathbf{A} - \mathfrak{p}\) 使 \(u_{f} \colon \mathbf{M}_{f} \to \mathbf{N}_{f}\) 是满射。
\item
  设 \(u_{\mathfrak{p}}\) 是双射，\(\mathbf{M}\) 是有限生成的且 \(\mathbf{N}\) 是有限表示的。则存在 \(f \in \mathbf{A} - \mathfrak{p}\) 使 \(u_{f}\) 是双射。
\end{enumerate}

设 R 与 Q 是 u 的核与余核；若 \(g \in A\)，则 \(u_g\)（相应地，\(u_p\)）的核与余核是 \(R_g\) 与 \(Q_g\)（相应地，\(R_p\) 与 \(Q_p\)）（\(\S 2\) 第 4 节定理 1）。则 \(Q_p = 0\)；由于 N 是有限生成的，Q 也是，且存在 \(g' \in A - p\) 使 \(g'Q = 0\)（\(\S 2\) 第 2 节命题 4 的推论 2），从而 \(Q_{g'} = 0\)。在 (ii) 的假设下，序列 \(0 \to R_{g'} \to M_{g'} \to N_{g'} \to 0\) 是正合的，因而 \(R_{g'}\) 是有限生成的（第 I 章 \(\S 2\) 第 8 节引理 9）。现在，

\[
\left(\mathrm{R} _ {g ^ {\prime}}\right) _ {\mathfrak{p}\mathbf{A}_{g ^ {\prime}}} = \mathrm{R} _ {\mathfrak{p}} = 0;
\]

因而存在 \(g_1 \in \mathbf{A}_{g'} - \mathfrak{p}\mathbf{A}_{g'}\) 使 \(g_1\mathbf{R}_{g'} = 0\)（§2 第 2 节命题 4 的推论 2）。于是 \(g_1 = g'' / g'^h\)，其中 \(g'' \in \mathbf{A} - \mathfrak{p}\)；由于 \(g' / 1\) 在 \(\mathbf{R}_{g'}\) 中可逆，故 \((g'' / 1)\mathbf{R}_{g'} = 0\)，从而 \(\mathbf{R}_{g'g''} = (\mathbf{R}_{g'})_{g'' / 1} = 0\)。若 \(f = g'g''\)，则 \(f \in \mathbf{A} - \mathfrak{p}\)、\(\mathbf{Q}_f = 0\) 且 \(\mathbf{R}_f = 0\)，于是 \(u_f\) 是双射。

推论. 若 N 是有限表示的且 \(N_{p}\) 是秩 p 的自由 \(A_{p}\) -模，则存在 \(f \in \mathbf{A} - \mathfrak{p}\) 使 \(N_{f}\) 是秩 p 的自由 \(A_{f}\) -模。

按假设存在 \(p\) 个元素 \(x_{i} \in \mathbf{N}\)（\(1 \leqslant i \leqslant p\)）使诸 \(x_{i}/1\) 构成自由 \(\mathbf{A}_{\mathfrak{p}}\) -模 \(\mathbf{N}_{\mathfrak{p}}\) 的一个基。考虑同态 \(u: \mathbf{A}^p \to \mathbf{N}\)（使 \(u(e_{i}) = x_{i}\)，对 \(1 \leqslant i \leqslant p\) 成立），其中 \((e_{i})_{1 \leqslant i \leqslant p}\) 是 \(\mathbf{A}^p\) 的典范基。由于按假设 \(u_{\mathfrak{p}}\) 是双射，凭命题 2 存在 \(f \in \mathbf{A} - \mathfrak{p}\) 使 \(u_{f}\) 是双射。

命题 3. 设 \((f_{i})_{i\in I}\) 是环 A 的元素的有限族，生成 A 的理想 A。则环 \(\mathbf{B} = \prod_{i\in I}\mathbf{A}_{f_i}\) 是一个忠实平坦 A-模。

由 §2 第 4 节定理 1，诸 \(\mathbf{A}_{f_i}\) 中的每一个都是平坦 A-模，因而 B 也是（第 I 章 §2 第 3 节命题 2）。另一方面，若 \(\mathfrak{p}\) 是 A 的一个素理想，则存在指标 \(i\) 使 \(f_i \notin \mathfrak{p}\)，于是 \(\mathfrak{p}_{f_i} = \mathfrak{p} \mathbf{A}_{f_i}\) 是 \(\mathbf{A}_{f_i}\) 的一个素理想。则 \(\mathfrak{p} \mathbf{B} \subset \mathfrak{p} \mathbf{A}_{f_i} \times \prod_{j \neq i} \mathbf{A}_{f_j} \neq \mathbf{B}\)，因为 \(\mathfrak{p} \mathbf{A}_{f_i} \neq \mathbf{A}_{f_i}\)；这足以蕴含 B 是忠实平坦 A-模（第 I 章 §3 第 1 节命题 1）。

推论. 在命题 3 的假设下，为使 A-模 M 是有限生成的（相应地，有限表示的），必须且只需对每个指标 i，\(A_{f_{i}}\) -模 \(M_{f_{i}}\) 是有限生成的（相应地，有限表示的）。

该条件显然是必要的（§2 第 4 节）。反之，若所有 \(M_{f_{i}}\) 都是有限生成的（相应地，有限表示的），则 \(M' = \prod_{i \in I} M_{f_{i}}\) 是有限生成的（相应地，有限表示的）B-模，因为我们显然可以假设对每个 i 有正合列 \(A_{f_i}^{m} \to A_{f_i}^{n} \to M_{f_i} \to 0\)，其中 m 与 n 与 i 无关。现在，\(M' = M \otimes_{A} B\)。该推论于是由命题 3 及第 I 章 §3 第 6 节命题 11 得出。

注意，关于诸 \(f_{i}\) 的条件表示开集 \(\mathbf{X}_{f_{i}}\) 构成 \(\operatorname{Spec}(\mathbf{A})\) 的一个覆盖（\(\S 4\) 第 3 节命题 11 的推论 3）。

